\documentclass[aspectratio=169]{beamer}

% Tema UFS — Universidade Federal de Sergipe
\usetheme{UFS}

% Pacotes base
\usepackage[utf8]{inputenc}
\usepackage[portuguese]{babel}
\usepackage{amsmath,amssymb}
\usepackage{graphicx}
\usepackage{booktabs}
\usepackage{xcolor}

% TikZ e bibliotecas
\usepackage{tikz}
\usetikzlibrary{shapes.geometric, arrows.meta, positioning, matrix, fit, calc,
  backgrounds, decorations.pathreplacing, shadows, patterns}

% Listagens — paleta VS Code Light+
\usepackage{listings}
\definecolor{bgcolor}{HTML}{FFFFFF}
\definecolor{linecolor}{HTML}{DDDDDD}
\definecolor{numcolor}{HTML}{999999}
\definecolor{kwcolor}{HTML}{0000FF}
\definecolor{strcolor}{HTML}{A31515}
\definecolor{cmtcolor}{HTML}{008000}

\lstdefinestyle{ufsStyle}{
  backgroundcolor=\color{bgcolor},
  basicstyle=\ttfamily\footnotesize\color{black},
  keywordstyle=\color{kwcolor}\bfseries,
  stringstyle=\color{strcolor},
  commentstyle=\color{cmtcolor}\itshape,
  numberstyle=\tiny\color{numcolor},
  numbers=left, numbersep=12pt,
  xleftmargin=20pt, framexleftmargin=20pt,
  breaklines=true, tabsize=4,
  keepspaces=true, showspaces=false,
  showstringspaces=false, showtabs=false,
  frame=single, rulecolor=\color{linecolor},
  framesep=4pt, captionpos=b, upquote=true,
  columns=fullflexible,
}
\lstset{style=ufsStyle, language=C}

% --- Metadados ---
\title{Algoritmos e Estruturas de Dados II}
\subtitle{Corretude de algoritmos iterativos e invariantes de laço --- Parte 2}
\author{Prof.\ André Yoshiaki Kashiwabara}
\institute{DCOMP -- Universidade Federal de Sergipe\\
\texttt{andreyoshiaki@dcomp.ufs.br}}
\date{}

\begin{document}

\begin{frame}[plain]
  \titlepage
\end{frame}

\begin{frame}{O que você vai aprender hoje}
  \begin{center}
  \begin{tikzpicture}[
    item/.style={rectangle, draw=UFSBlue, fill=UFSBlue!10, rounded corners,
      minimum width=9.5cm, minimum height=0.75cm, font=\small}
  ]
    \node[item] (t1) at (0,0)    {Relembrar o método: invariante $+$ inicialização, manutenção e término};
    \node[item] (t2) at (0,-1.1) {Prova de corretude da \textbf{busca binária} --- e por que tantas implementações erram};
    \node[item] (t3) at (0,-2.2) {Prova de corretude do \textbf{particionamento de Lomuto} (coração do quicksort)};
    \node[item] (t4) at (0,-3.3) {Usar o invariante como \textbf{ferramenta de projeto}, não só de prova};
    \draw[->, thick, UFSBlue] (t1)--(t2);
    \draw[->, thick, UFSBlue] (t2)--(t3);
    \draw[->, thick, UFSBlue] (t3)--(t4);
  \end{tikzpicture}
  \end{center}
\end{frame}

\begin{frame}{Agenda}
  \tableofcontents
\end{frame}

% =====================================================================
\section{Relembrando o método}
% =====================================================================

\begin{frame}{O método em uma figura}
  \begin{center}
  \begin{tikzpicture}[
    fase/.style={rectangle, draw=UFSBlue, fill=UFSBlue!10, rounded corners,
      minimum width=3.6cm, minimum height=1.2cm, font=\small, align=center},
    seta/.style={->, thick, UFSBlue}
  ]
    \node[fase] (ini) at (0,0)
      {\textbf{Inicialização}\\ vale \emph{antes} da\\ primeira iteração};
    \node[fase] (man) at (5,0)
      {\textbf{Manutenção}\\ cada iteração\\ preserva o invariante};
    \node[fase] (ter) at (10,0)
      {\textbf{Término}\\ ao sair, invariante\\ $\Rightarrow$ poscondição};
    \draw[seta] (ini)--(man);
    \draw[seta] (man)--(ter);
    \draw[seta] (man) to[out=140, in=40, looseness=4] (man);
  \end{tikzpicture}
  \end{center}
  \begin{itemize}
    \item Na Parte 1: busca do máximo (invariante sobre um valor) e ordenação por
      inserção (invariante estrutural, laços aninhados).
    \item Hoje: busca binária e particionamento --- e, para corretude total, uma medida
      inteira não negativa que decresce a cada iteração (garantia de término).
  \end{itemize}
\end{frame}

% =====================================================================
\section{Terceiro algoritmo: busca binária}
% =====================================================================

\begin{frame}{Por que a busca binária merece uma prova?}
  \begin{columns}[T]
    \begin{column}{0.5\textwidth}
      \textbf{Problema:}\\[2pt]
      Dado $A[0..n-1]$ \textbf{ordenado} e um valor $x$, devolver um índice $i$ com
      $A[i] = x$, ou $-1$ se $x$ não está no vetor.
      \begin{itemize}
        \item Publicada em 1946; a primeira versão correta para todo $n$ demorou
          mais de uma década a aparecer.
        \item Bentley: a maioria dos programadores profissionais não a implementa
          corretamente de primeira.
      \end{itemize}
    \end{column}
    \begin{column}{0.5\textwidth}
      \textbf{Solução:}\\[2pt]
      Manter uma \emph{janela} $[lo..hi]$ onde $x$ \emph{ainda pode} estar e cortá-la
      pela metade a cada comparação.
      \begin{itemize}
        \item O invariante diz exatamente o que a janela significa --- e dita como
          atualizar \texttt{lo} e \texttt{hi} sem erro.
      \end{itemize}
    \end{column}
  \end{columns}
\end{frame}

\begin{frame}[fragile]{Busca binária: o algoritmo}
\begin{lstlisting}
int busca_binaria(int A[], int n, int x) { /* pre: A ordenado */
    int lo = 0, hi = n - 1;
    while (lo <= hi) {
        /* invariante: se x esta em A, esta em A[lo..hi] */
        int meio = lo + (hi - lo) / 2;
        if (A[meio] == x)      return meio;
        else if (A[meio] < x)  lo = meio + 1;
        else                   hi = meio - 1;
    }
    return -1;   /* janela vazia: x nao esta no vetor    */
}
\end{lstlisting}
  \begin{block}{Invariante}
    No início de cada iteração: \textbf{se} $x$ ocorre em $A[0..n-1]$,
    \textbf{então} $x$ ocorre em $A[lo..hi]$.
  \end{block}
\end{frame}

\begin{frame}{Busca binária: o invariante em figura}
  \begin{center}
  \begin{tikzpicture}[
    cel/.style={draw, minimum width=0.95cm, minimum height=0.8cm, font=\small},
    fora/.style={cel, fill=UFSLightGray},
    dentro/.style={cel, fill=UFSBlue!20},
    meio/.style={cel, fill=UFSYellow!50}
  ]
    \node[fora]   (a0) at (0,0)    {1};
    \node[fora]   (a1) at (0.95,0) {3};
    \node[dentro] (a2) at (1.9,0)  {4};
    \node[dentro] (a3) at (2.85,0) {7};
    \node[meio]   (a4) at (3.8,0)  {9};
    \node[dentro] (a5) at (4.75,0) {11};
    \node[dentro] (a6) at (5.7,0)  {15};
    \node[fora]   (a7) at (6.65,0) {20};

    \foreach \i/\n in {0/a0,1/a1,2/a2,3/a3,4/a4,5/a5,6/a6,7/a7}
      \node[below=1pt of \n, font=\scriptsize, color=UFSGray] {\i};

    \node[below=14pt of a2, font=\small, color=UFSBlue] {$lo$};
    \node[below=14pt of a6, font=\small, color=UFSBlue] {$hi$};
    \node[below=14pt of a4, font=\small] {$meio$};

    \draw[decorate, decoration={brace, amplitude=5pt}, thick, UFSBlue]
      (1.43,0.5) -- (6.17,0.5)
      node[midway, above=6pt, font=\small, color=UFSBlue]
      {se $x$ está no vetor, está aqui};
  \end{tikzpicture}
  \end{center}
  \begin{itemize}
    \item Cinza: regiões já \textbf{descartadas} --- a ordenação garante que $x$ não está nelas.
    \item O invariante é \emph{condicional}: não promete que $x$ está na janela; promete
      que \textbf{fora dela é impossível}.
    \item Janela vazia ($lo > hi$) $+$ invariante $\Rightarrow$ $x$ não está no vetor: o
      \texttt{return -1} está \emph{provado}, não ``chutado''.
  \end{itemize}
\end{frame}

\begin{frame}{Busca binária: prova de corretude}
  \begin{itemize}
    \item \textbf{Inicialização}: $[lo..hi] = [0..n-1]$ é o vetor inteiro ---
      a implicação vale trivialmente. \checkmark
    \item \textbf{Manutenção}: seja $m$ o meio. Se $A[m] < x$, a ordenação dá
      $A[k] \leq A[m] < x$ para todo $k \leq m$: descartar $[lo..m]$ é seguro, e
      $lo = m + 1$ preserva o invariante. Caso $A[m] > x$ é simétrico. \checkmark
    \item \textbf{Término com retorno no laço}: só retornamos $meio$ quando
      $A[meio] = x$ --- correto por inspeção. \checkmark
    \item \textbf{Término com janela vazia}: $lo > hi$ e o invariante
      (contrapositiva): $x$ não está em $A[lo..hi]$ vazio, logo não está no vetor
      --- devolver $-1$ é correto. \checkmark
    \item \textbf{O laço termina?} A medida $hi - lo + 1$ (tamanho da janela) é
      $\geq 0$ e \textbf{decresce estritamente}: $meio$ está em $[lo..hi]$, e os dois
      ramos excluem $meio$ da nova janela. \checkmark
  \end{itemize}
\end{frame}

\begin{frame}[fragile]{Exemplo: dois erros clássicos que a prova denuncia}
  \begin{exampleblock}{Erro 1: \texttt{lo = meio} em vez de \texttt{lo = meio + 1}}
    Com janela $[3..4]$, $meio = 3$ e $A[3] < x$: a ``nova'' janela continua $[3..4]$.
    A medida $hi - lo + 1$ \textbf{não decresce} --- laço infinito. A prova de término
    exige excluir $meio$.
  \end{exampleblock}
  \begin{exampleblock}{Erro 2: \texttt{meio = (lo + hi) / 2}}
    Matematicamente idêntico --- mas em C, com $n$ grande, $lo + hi$ pode estourar o
    \texttt{int} (overflow) e tornar \texttt{meio} negativo.
    \texttt{lo + (hi - lo) / 2} calcula o mesmo valor sem estourar.
  \end{exampleblock}
  \begin{itemize}
    \item O primeiro erro é \emph{lógico}: a prova o captura no passo de término.
    \item O segundo é da \emph{máquina}: a prova vale sobre inteiros ideais; a
      implementação precisa respeitar os limites do tipo.
  \end{itemize}
\end{frame}

\begin{frame}{Recapitulando: busca binária}
  \begin{block}{O que aprendemos}
    \begin{itemize}
      \item Invariante condicional: ``se $x$ está no vetor, está na janela $[lo..hi]$''.
      \item A manutenção usa a \textbf{precondição} (vetor ordenado) para justificar o descarte.
      \item Término via medida decrescente: o tamanho da janela cai a cada iteração ---
        e é isso que proíbe \texttt{lo = meio}.
      \item O caso ``não encontrado'' sai \emph{provado} do invariante, de graça.
    \end{itemize}
  \end{block}
\end{frame}

% =====================================================================
\section{Quarto algoritmo: particionamento de Lomuto}
% =====================================================================

\begin{frame}{Por que estudar o particionamento?}
  \begin{columns}[T]
    \begin{column}{0.5\textwidth}
      \textbf{Problema:}\\[2pt]
      Reorganizar $A[p..r]$ em torno do pivô $A[r]$: menores ou iguais à esquerda,
      maiores à direita, pivô no meio --- devolvendo a posição final do pivô.
      \begin{itemize}
        \item É o coração do \textbf{quicksort}: se o particionamento está certo,
          a recursão faz o resto.
      \end{itemize}
    \end{column}
    \begin{column}{0.5\textwidth}
      \textbf{Solução:}\\[2pt]
      Varrer o trecho mantendo \textbf{quatro regiões} bem delimitadas por dois
      índices ($i$ e $j$).
      \begin{itemize}
        \item O invariante aqui é uma \emph{fotografia} das quatro regiões ---
          o exemplo mais visual dos quatro algoritmos que estudamos.
      \end{itemize}
    \end{column}
  \end{columns}
\end{frame}

\begin{frame}[fragile]{Particionamento de Lomuto: o algoritmo}
\begin{lstlisting}
int particiona(int A[], int p, int r) {
    int pivo = A[r];
    int i = p - 1;
    for (int j = p; j < r; j++) {
        /* invariante: A[p..i] <= pivo,                  */
        /*   A[i+1..j-1] > pivo, A[j..r-1] nao examinado */
        if (A[j] <= pivo) {
            i++;
            troca(&A[i], &A[j]);
        }
    }
    troca(&A[i + 1], &A[r]);
    return i + 1;
}
\end{lstlisting}
\end{frame}

\begin{frame}{O invariante das quatro regiões}
  \begin{center}
  \begin{tikzpicture}[
    cel/.style={draw, minimum width=0.95cm, minimum height=0.8cm, font=\small},
    menor/.style={cel, fill=UFSBlue!20},
    maior/.style={cel, fill=UFSRed!20},
    resto/.style={cel, fill=UFSLightGray},
    piv/.style={cel, fill=UFSYellow!50}
  ]
    \node[menor] (a0) at (0,0)    {2};
    \node[menor] (a1) at (0.95,0) {1};
    \node[maior] (a2) at (1.9,0)  {8};
    \node[maior] (a3) at (2.85,0) {7};
    \node[maior] (a4) at (3.8,0)  {9};
    \node[resto] (a5) at (4.75,0) {3};
    \node[resto] (a6) at (5.7,0)  {6};
    \node[piv]   (a7) at (6.65,0) {5};

    \foreach \i/\n in {0/a0,1/a1,2/a2,3/a3,4/a4,5/a5,6/a6,7/a7}
      \node[below=1pt of \n, font=\scriptsize, color=UFSGray] {\i};

    \node[below=14pt of a1, font=\small, color=UFSBlue] {$i$};
    \node[below=14pt of a5, font=\small, color=UFSGray] {$j$};
    \node[above=6pt of a7, font=\small, color=UFSGray] {pivô};

    \draw[decorate, decoration={brace, amplitude=4pt}, thick, UFSBlue]
      (-0.47,0.5) -- (1.42,0.5)
      node[midway, above=5pt, font=\scriptsize, color=UFSBlue] {$\leq 5$};
    \draw[decorate, decoration={brace, amplitude=4pt}, thick, UFSRed]
      (1.43,0.5) -- (4.27,0.5)
      node[midway, above=5pt, font=\scriptsize, color=UFSRed] {$> 5$};
    \draw[decorate, decoration={brace, amplitude=4pt}, thick, UFSGray]
      (4.28,0.5) -- (6.17,0.5)
      node[midway, above=5pt, font=\scriptsize, color=UFSGray] {não examinado};
  \end{tikzpicture}
  \end{center}
  \begin{block}{Invariante (início da iteração $j$)}
    $A[p..i] \leq \text{pivô}$ \quad e \quad $A[i+1..j-1] > \text{pivô}$ \quad e
    \quad $A[j..r-1]$ não examinado \quad e \quad $A[r] = \text{pivô}$.
  \end{block}
\end{frame}

\begin{frame}{Manutenção: os dois casos da iteração}
  \begin{center}
  \begin{tikzpicture}[
    cel/.style={draw, minimum width=0.75cm, minimum height=0.65cm, font=\scriptsize},
    menor/.style={cel, fill=UFSBlue!20},
    maior/.style={cel, fill=UFSRed!20},
    resto/.style={cel, fill=UFSLightGray},
    rot/.style={font=\small, anchor=west}
  ]
    % Caso 1
    \node[rot] at (-1.0,0.9) {\textbf{Caso $A[j] > $ pivô}: só avançar $j$ --- a região vermelha cresce.};
    \node[menor] at (0,0)    {$\leq$};
    \node[menor] at (0.75,0) {$\leq$};
    \node[maior] at (1.5,0)  {$>$};
    \node[maior] at (2.25,0) {$>$};
    \node[maior, very thick, draw=UFSRed] at (3.0,0) {$>$};
    \node[resto] at (3.75,0) {?};
    \node[font=\scriptsize, color=UFSGray] at (3.0,-0.6) {$j$};

    % Caso 2
    \node[rot] at (-1.0,-1.5) {\textbf{Caso $A[j] \leq$ pivô}: \texttt{i++} e trocar $A[i] \leftrightarrow A[j]$ --- a azul cresce.};
    \node[menor] at (0,-2.3)    {$\leq$};
    \node[menor] at (0.75,-2.3) {$\leq$};
    \node[menor, very thick, draw=UFSBlue] at (1.5,-2.3) {$\leq$};
    \node[maior] at (2.25,-2.3) {$>$};
    \node[maior] at (3.0,-2.3)  {$>$};
    \node[resto] at (3.75,-2.3) {?};
    \node[font=\scriptsize, color=UFSGray] at (1.2,-2.95) {$i$};
    \draw[->, thick, UFSBlue] (3.0,-2.65) to[out=-60, in=-120] (1.5,-2.65);
  \end{tikzpicture}
  \end{center}
  \begin{itemize}
    \item A troca move o \emph{primeiro} elemento da região vermelha para o fim dela ---
      as duas regiões continuam homogêneas; se a vermelha está vazia ($i + 1 = j$),
      a ``troca'' é de $A[j]$ com ele mesmo e o invariante segue valendo.
  \end{itemize}
\end{frame}

\begin{frame}{Particionamento: prova de corretude}
  \begin{itemize}
    \item \textbf{Inicialização} ($j = p$, $i = p - 1$): as regiões $A[p..p-1]$ e
      $A[p..p-1]$ são vazias --- as duas cláusulas valem por vacuidade. \checkmark
    \item \textbf{Manutenção}: os dois casos do slide anterior preservam as quatro
      regiões. \checkmark
    \item \textbf{Término}: com $j = r$, todo $A[p..r-1]$ está classificado:
      $A[p..i] \leq \text{pivô}$ e $A[i+1..r-1] > \text{pivô}$.
      A troca final põe o pivô em $A[i+1]$, entre as duas regiões ---
      exatamente a poscondição, com retorno $i + 1$. \checkmark
    \item \textbf{O laço termina?} É um \texttt{for} com $j$ de $p$ a $r - 1$:
      a medida $r - j$ decresce a cada iteração. \checkmark
    \item \textbf{Permutação:} só fazemos trocas --- nenhum elemento é criado ou
      perdido (mesma cláusula que na ordenação por inserção).
  \end{itemize}
\end{frame}

\begin{frame}{Exemplo: rastreando $A = [3, 8, 1, 5]$ com pivô $5$}
  \begin{center}
  \small
  \begin{tabular}{c c c l}
    \toprule
    Início de $j$ & $i$ & Vetor & Regiões $[\leq \mid > \mid ?]$ \\
    \midrule
    $j = 0$ & $-1$ & $[3, 8, 1 \mid 5]$ & $[\ \mid\ \mid 3, 8, 1]$ \\
    $j = 1$ & $0$  & $[\mathbf{3}, 8, 1 \mid 5]$ & $[3 \mid\ \mid 8, 1]$ \\
    $j = 2$ & $0$  & $[\mathbf{3}, \mathbf{8}, 1 \mid 5]$ & $[3 \mid 8 \mid 1]$ \\
    saída ($j = 3$) & $1$ & $[\mathbf{3}, \mathbf{1}, \mathbf{8} \mid 5]$ & $[3, 1 \mid 8 \mid\ ]$ \\
    troca final & --- & $[3, 1, \mathbf{5}, 8]$ & pivô na posição $2$ \checkmark \\
    \bottomrule
  \end{tabular}
  \end{center}
  \begin{itemize}
    \item Em $j = 2$: $A[2] = 1 \leq 5$, então $i$ vira $1$ e trocamos $A[1] \leftrightarrow A[2]$.
    \item A troca final coloca o pivô entre as regiões: tudo à esquerda $\leq 5$,
      tudo à direita $> 5$.
  \end{itemize}
\end{frame}

\begin{frame}{Recapitulando: particionamento de Lomuto}
  \begin{block}{O que aprendemos}
    \begin{itemize}
      \item Invariante ``geográfico'': o vetor dividido em quatro regiões com
        significados distintos.
      \item Inicialização por \textbf{vacuidade}: regiões vazias satisfazem qualquer
        propriedade.
      \item A manutenção é uma análise de dois casos, cada um fazendo uma região crescer.
      \item Com o particionamento provado, a corretude do quicksort segue por indução
        no tamanho do trecho.
    \end{itemize}
  \end{block}
\end{frame}

% =====================================================================
\section{Invariante como ferramenta de projeto}
% =====================================================================

\begin{frame}{Invariante primeiro, código depois}
  \begin{block}{A inversão de perspectiva}
    Em vez de escrever o laço e \emph{depois} procurar o invariante, escolha o
    invariante e deixe-o \textbf{ditar} o código.
  \end{block}
  \begin{itemize}
    \item \textbf{Inicialização} $\rightarrow$ diz como declarar e inicializar as variáveis.
    \item \textbf{Manutenção} $\rightarrow$ diz o que o corpo do laço \emph{precisa} fazer.
    \item \textbf{Término $+$ poscondição} $\rightarrow$ diz qual é a condição de parada.
  \end{itemize}
  \begin{exampleblock}{Nos quatro algoritmos que estudamos}
    Máximo: ``maior já visto''. \quad Inserção: ``prefixo ordenado''.\\
    Busca binária: ``$x$ só pode estar em $[lo..hi]$''. \quad Lomuto: ``quatro regiões''.
  \end{exampleblock}
\end{frame}

\begin{frame}{Resumo das duas aulas}
  \begin{center}
  \footnotesize
  \begin{tabular}{l l l}
    \toprule
    Algoritmo & Tipo de invariante & O que a prova exigiu \\
    \midrule
    Busca do máximo & valor acumulado & caso base não trivial (\texttt{max = A[0]}) \\
    Ordenação por inserção & prefixo estruturado & ordem $+$ permutação; laços aninhados \\
    Busca binária & janela condicional & precondição (ordenado) $+$ medida decrescente \\
    Particionamento & regiões geográficas & vacuidade $+$ análise de casos \\
    \bottomrule
  \end{tabular}
  \end{center}
  \begin{itemize}
    \item O método é sempre o mesmo: \textbf{inicialização, manutenção, término} ---
      o que muda é a criatividade para \emph{enunciar} o invariante.
    \item Enunciar o invariante certo é a metade difícil --- e é uma habilidade que se
      treina: o tutorial desta aula é esse treino.
  \end{itemize}
\end{frame}

\section*{Referências}
\begin{frame}[allowframebreaks]{Referências}
  \nocite{*}
  \bibliographystyle{plain}
  \bibliography{../referencias}
\end{frame}

\end{document}
